% Generated draft: requires mathematical and source review.
\documentclass[11pt]{article}
\usepackage[margin=25mm]{geometry}
\usepackage{fontspec}
\setmainfont{Latin Modern Roman}
\usepackage{amsmath,amssymb,mathtools,unicode-math}
\setmathfont{Latin Modern Math}
\usepackage{enumitem,needspace,xurl,hyperref}
\setlength{\parindent}{0pt}
\setlength{\parskip}{5pt}
\setlength{\emergencystretch}{3em}
\newcommand{\N}{\mathbb{N}}
\newcommand{\Z}{\mathbb{Z}}
\newcommand{\Q}{\mathbb{Q}}
\newcommand{\R}{\mathbb{R}}
\newcommand{\C}{\mathbb{C}}
\newcommand{\F}{\mathbb{F}}
\newcommand{\A}{\mathbb{A}}
\newcommand{\PP}{\mathbb{P}}
\newcommand{\E}{\mathbb{E}}
\newcommand{\eps}{\varepsilon}
\newcommand{\dd}{\,\mathrm{d}}
\begin{document}
\section*{UnsolvedMath ID: 7600012}
\textbf{Source problem code:} catalogue entry
\subsection{Definitions and mathematical statement}
The source asks whether some first-order theory has a minimum universality spectrum, but does not define the spectrum or its ordering.

One common universality setup fixes a theory $T$, an infinite cardinal $\kappa$, and an embedding notion. A $\kappa$-sized model $U\models T$ is universal if every $\kappa$-sized $M\models T$ embeds into $U$ in that sense. With elementary embeddings as an indicated convention, a candidate spectrum is
\[
\operatorname{Univ}(T)=\{\kappa:\ \exists U\models T,\ |U|=\kappa,
\ \forall M\models T\ (|M|=\kappa\Rightarrow M\preccurlyeq U\text{ up to an embedding})\}.
\]
Here the embedding notation means that $M$ is elementarily isomorphic to a substructure, not that the given underlying set of $M$ is literally contained in $U$.

If minimum means inclusion-minimum among such spectra, the target would be
\[
\exists T_*\ \forall T,\qquad\operatorname{Univ}(T_*)\subseteq\operatorname{Univ}(T).
\]
This is only a possible interpretation. A spectrum may instead encode least sizes of jointly universal families, or an associated comparison invariant. A least element under an order is also stronger than merely a minimal element with no strictly smaller competitor.

The record specifies neither the admissible theories (complete, consistent, infinite-model, countable-language, or another class), the embedding type, cardinal range, set-theoretic background, nor the comparison order. Allowing theories with only finite models would trivialize the displayed infinite-cardinal spectrum by making it empty, showing why these restrictions matter. The source question therefore cannot be assigned the displayed inclusion interpretation as a unique precise theorem without its original definitions.
\subsection{Short English statement}
Does the intended class of first-order theories contain a least universality spectrum? The spectrum, comparison order, theory class and embedding convention are missing, and minimal versus least must be distinguished.
\subsection{Sources}
\begin{itemize}
\item Local UnsolvedMath registry, permanent ID 7600012, exact source code catalogue entry, ``The universality spectrum problem''; no proposer or year supplied. Registry status open, classification OPEN-TRIAGE.
\item Wikipedia contributors, \emph{List of unsolved problems in mathematics}, Model theory and formal languages, bullet 12, pinned revision 1366636767, \url{https://en.wikipedia.org/w/index.php?title=List_of_unsolved_problems_in_mathematics&oldid=1366636767#Model_theory_and_formal_languages}. Supplied MediaWiki extraction accessed July 29, 2026. No primary definition or citation is supplied.
\item The research summary associates the problem with Shelah and ``Malick'' and unspecified thesis/papers, but expressly provides no verified decisive citation. Those attributions and the relevant notion of spectrum require checking; no theory-comparison theorem is inferred.
\item Provenance flags primary-source status NEEDS\_REVIEW and says the bullet remains unpublished pending review despite published=true metadata. Wikipedia text is CC BY-SA 4.0; attribution and share-alike handling remain required for release.
\end{itemize}
\end{document}
